{\small\textbf{Resilience gained (Zheng, 1:00) [21:00–22:00]}}\par\smallskip
Third, what the Netherlands gained.
\begin{itemize}\setlength\itemsep{0pt}
\item In 2020 it had \textbf{no} capital that was releasable by design. Today it has about \textbf{6.7 billion euros}, roughly half of the peak losses after the financial crisis, enough to support up to 150 billion of lending.
\item ECB research shows why this is a good deal: raising requirements in normal times costs very little lending, about 0.1 percent per point, while releasing in bad times can boost lending by up to ten percent.
\end{itemize}
\par\smallskip
And its role is growing. The mortgage risk-weight floor expires at the end of November 2026, and DNB says this "underpins the importance" of the two-percent buffer.
